\chapter[Lapisan Teknologi: Node, Artefak, dan Deployment]{Lapisan Teknologi:
Node, Artefak, dan Deployment Sistem Transfer Dana}
\label{ch:technology-architecture}
\setlength{\emergencystretch}{2em}

\section*{Tujuan Pembelajaran}
\begin{itemize}
  \item Menjelaskan elemen lapisan teknologi ArchiMate---node, perangkat
        (\emph{device}), perangkat lunak sistem (\emph{system software}), layanan
        teknologi, artefak (\emph{artifact}), dan jaringan komunikasi.
  \item Menyusun arsitektur teknologi acuan tempat komponen aplikasi di-deploy
        (lewat artefak) dan objek data diwujudkan pada infrastruktur.
  \item Menempatkan analisis teknologi dalam TOGAF ADM \textbf{Fase D (Arsitektur
        Teknologi)}: menyusun \emph{baseline}, \emph{target}, dan analisis
        kesenjangan.
  \item Menutup rantai keterlacakan motivasi--strategi--bisnis--aplikasi--data--
        teknologi, serta mengaitkannya dengan pola deployment dan konteks
        infrastruktur Indonesia.
\end{itemize}

Tujuan pertama dan kedua mendukung CLO 3; tujuan ketiga mendukung CLO 2 dan
CLO 3; tujuan keempat mendukung CLO 6 dan CLO 7 karena mahasiswa menautkan
rancangan teknologi dengan regulasi infrastruktur dan argumen akademik.

\section{Ringkasan Awal Bab}

Bab~\ref{ch:application-architecture} memetakan \emph{perangkat lunak apa} yang
menjalankan proses dan mengolah data. Bab ini menutup lapisan teknis dengan
menjawab \emph{di atas apa} perangkat lunak itu berjalan: node, perangkat,
perangkat lunak sistem, jaringan, dan artefak yang tersimpan. Dalam ArchiMate,
lapisan teknologi memuat \textbf{node} (sumber daya komputasi), \textbf{device}
(perangkat fisik), \textbf{system software} (sistem operasi, runtime, DBMS),
\textbf{technology service}, \textbf{communication network}, dan
\textbf{artifact} (wujud fisik data dan komponen).

Di sinilah rantai keterlacakan tertutup. \textbf{Node men-deploy artefak}, dan
\textbf{artefak itu mewujudkan} komponen aplikasi
(Bab~\ref{ch:application-architecture}) maupun objek data
(Bab~\ref{ch:data-architecture}). Dengan demikian setiap kebutuhan dapat
ditelusuri dari motivasi sampai ke node dan artefak yang menjalankannya.
Keluaran bab memperkuat bagian \textbf{Analisis} dan menjadi \textbf{kerangka
akhir} naskah.

\section{Pendahuluan}

Teknologi sering dianggap urusan paling teknis, padahal keputusannya sangat
manajerial: pusat data sendiri atau awan (\emph{cloud}), satu server besar atau
banyak node kecil, mengganti atau membungkus sistem inti lama. Keputusan ini
menentukan biaya, keandalan, kepatuhan, dan kelincahan. Arsitektur teknologi yang
baik membuat pilihan itu terlihat---node apa menjalankan komponen apa, artefak
apa menyimpan data apa, dan jaringan apa menghubungkannya---tanpa perlu masuk ke
konfigurasi rinci.

Bagi makalah, lapisan teknologi mengubah rekomendasi menjadi konkret dan dapat
dinilai kelayakannya. Pilihan kontainer dan deployment bertahap dapat
dijustifikasi dengan pembahasan modularisasi dan deployment perangkat lunak
modern \cite{newman2021building}, sedangkan kewajiban penyelenggaraan sistem
elektronik memberi batas dan konteks \cite{pp712019}.

\section{Skenario Pembuka dan Pertanyaan Pemandu}

Setelah komponen aplikasi rel transfer dirancang, tim infrastruktur Bank Wanua
menghadapi pertanyaan: di mana komponen di-deploy? Sistem inti lama berjalan di
\emph{mainframe} yang tidak dapat dipindah cepat, sedangkan komponen baru cocok
di-deploy sebagai kontainer. Koneksi ke BI-FAST menuntut jaringan yang aman dan
andal, dan data transaksi menuntut basis data yang tahan gagal serta berjejak.

\subsection*{Pertanyaan Pemandu}
\begin{enumerate}
  \item Node dan perangkat lunak sistem apa yang diperlukan untuk men-deploy
        komponen aplikasi?
  \item Artefak apa yang mewujudkan objek data, dan di node mana ia disimpan?
  \item Layanan teknologi dan jaringan apa yang menopang lanskap?
  \item Bagaimana sistem inti lama diintegrasikan tanpa penggantian sekaligus?
  \item Bagaimana kondisi infrastruktur saat ini berbeda dari target, dan apa
        kesenjangannya?
\end{enumerate}

\section{Mengapa Bab Ini Penting bagi Pembaca Non-TI}

Pembaca non-TI tidak perlu mengelola server, tetapi perlu memahami keputusan
teknologi yang menentukan biaya dan risiko: awan atau pusat data sendiri,
kedaulatan data, ketahanan terhadap gangguan, dan kepatuhan. Dengan beberapa
notasi sederhana---node, perangkat lunak sistem, layanan teknologi, artefak, dan
jaringan---seorang analis dapat menjelaskan arsitektur teknologi dan trade-off-nya
sebagai keputusan manajemen, bukan sekadar konfigurasi.

\section{Konsep Inti dan Glosarium Mini}

\subsection{Lapisan Teknologi ArchiMate}

\textbf{Node} adalah sumber daya komputasi yang menjalankan perilaku dan
men-deploy perangkat lunak, misalnya kluster kontainer. \textbf{Device} adalah
node fisik, misalnya server di pusat data. \textbf{System software} adalah
perangkat lunak sistem, seperti sistem operasi, runtime kontainer, atau DBMS.
\textbf{Technology service} adalah layanan yang ditawarkan node, misalnya layanan
komputasi. \textbf{Communication network} menghubungkan node. \textbf{Artifact}
adalah wujud fisik data atau perangkat lunak yang dapat di-deploy---misalnya
\emph{container image} atau skema basis data.

\subsection{Notasi Elemen Teknologi di Archi}

Gambar~\ref{fig:technology-legend} menampilkan jenis elemen lapisan teknologi di
Archi. Semua elemen teknologi berwarna kehijauan; ikon di pojok kanan atas
membedakan jenisnya.

\begin{figure}[htbp]
\centering
\includegraphics[width=0.92\textwidth]{technology-legend.pdf}
\caption{Jenis elemen lapisan teknologi ArchiMate dan simbolnya di Archi.}
\label{fig:technology-legend}
\end{figure}

\begin{table}[htbp]
\centering
\scriptsize
\caption{Elemen lapisan teknologi ArchiMate, kegunaan, dan contohnya.}
\label{tab:tech-elements}
\begin{tabularx}{\textwidth}{@{}p{0.24\textwidth} p{0.22\textwidth} X@{}}
\hline
\textbf{Elemen} & \textbf{Kegunaan} & \textbf{Contoh} \\
\hline
Node & Sumber daya komputasi yang men-deploy perangkat lunak. & Kluster Kontainer,
Server Basis Data, Sistem Inti (Mainframe). \\
Device & Node fisik (perangkat keras). & Server Fisik (Data Center). \\
System Software & Sistem operasi, runtime, atau DBMS. & Runtime Kontainer \& OS,
Sistem Basis Data (PostgreSQL). \\
Technology Service & Layanan yang ditawarkan node ke lapisan aplikasi. & Layanan
Komputasi Kontainer, Layanan Basis Data. \\
Artifact & Wujud fisik data/komponen yang dapat di-deploy. & Container Image, Skema
Basis Data, Berkas Ledger. \\
Communication Network & Penghubung antar-node. & Jaringan Internal Bank, Koneksi
Aman BI-FAST. \\
Path & Jalur logis komunikasi antar-node. & (mis.) jalur antar pusat data. \\
\hline
\end{tabularx}
\end{table}

\subsection{Deployment: Node, Artefak, dan Komponen}

Cara sebuah komponen aplikasi ``sampai'' ke infrastruktur disebut deployment
(\emph{deployment}). Penting dipahami: dalam ArchiMate sebuah node \textbf{tidak}
langsung terhubung ke komponen aplikasi. Sesuai spesifikasi The Open Group,
``sebuah komponen aplikasi diwujudkan oleh sebuah artefak dan, secara tidak
langsung, oleh node tempat artefak itu di-deploy''. Deployment karena itu memakai
dua relasi berurutan (Gambar~\ref{fig:deployment}): sebuah \textbf{node
men-deploy} (relasi \emph{assignment}) sebuah \textbf{artefak}, dan artefak itu
\textbf{mewujudkan} (relasi \emph{realization}) komponen aplikasi maupun objek
data. Node tidak ``menjadi'' komponen; ia menyediakan tempat artefak di-deploy,
sedangkan artefak adalah bentuk fisik yang dapat disimpan dan disalin. Sebagai
mekanisme ketiga, sebuah \textbf{layanan teknologi} (yang diwujudkan node)
\textbf{melayani} komponen atau fungsi aplikasi---inilah satu-satunya relasi
\emph{serving} yang sah dari teknologi ke aplikasi (node tidak melayani komponen
secara langsung).

\begin{figure}[htbp]
\centering
\scalebox{0.95}{
\begin{tikzpicture}[
  font=\footnotesize,
  app/.style={rectangle, rounded corners=1.5mm, draw=blue!70!black, fill=blue!12,
    thick, align=center, text width=34mm, minimum height=11mm},
  tech/.style={rectangle, rounded corners=1.5mm, draw=praditagreen,
    fill=praditagreen!16, thick, align=center, text width=34mm, minimum height=11mm},
  data/.style={rectangle, rounded corners=1.5mm, draw=cyan!60!black, fill=cyan!12,
    thick, align=center, text width=30mm, minimum height=11mm},
  asg/.style={-Latex, thick, draw=gray!70},
  real/.style={-{Triangle[open]}, thick, draw=gray!70, densely dashed}
]
  \node[app] (comp) {Komponen Aplikasi};
  \node[tech, below=13mm of comp] (art) {Artifact};
  \node[data, below=13mm of art] (do) {Objek Data};
  \node[tech, left=22mm of art] (node) {Node};
  \draw[asg] (node) -- node[above,font=\scriptsize]{men-deploy} (art);
  \draw[real] (art) -- node[right,font=\scriptsize]{mewujudkan} (comp);
  \draw[real] (art) -- node[right,font=\scriptsize]{mewujudkan} (do);
\end{tikzpicture}
}
\caption{Pola deployment (kanonik ArchiMate): node \emph{men-deploy} (assignment)
artefak; artefak \emph{mewujudkan} (realization) komponen aplikasi dan objek data.
Node menjalankan komponen secara tidak langsung, lewat artefak.}
\label{fig:deployment}
\end{figure}

\subsection{TOGAF ADM Fase D: Arsitektur Teknologi}

Fase D menyusun arsitektur teknologi acuan: node, platform, jaringan, dan pola
deployment---baseline dan target---lalu menganalisis kesenjangan \cite{togaf10}.
Keluaran utamanya adalah katalog teknologi, diagram deployment, dan matriks
komponen-terhadap-teknologi. Tabel~\ref{tab:adm-phase-d} memetakan langkahnya ke
kasus Bank Wanua.

\begin{table}[htbp]
\centering
\scriptsize
\caption{Pemetaan TOGAF ADM Fase D (Arsitektur Teknologi) ke artefak Bank Wanua.}
\label{tab:adm-phase-d}
\begin{tabularx}{\textwidth}{@{}p{0.27\textwidth} X@{}}
\hline
\textbf{Langkah Fase D} & \textbf{Artefak / keluaran pada kasus Bank Wanua} \\
\hline
Ruang lingkup dan prinsip teknologi & Lingkup infrastruktur rel transfer; prinsip
kontainer, ketahanan, keamanan, dan kedaulatan data. \\
Arsitektur teknologi baseline & Sistem inti di \emph{mainframe}; server terpisah;
integrasi dan jaringan belum seragam. \\
Arsitektur teknologi target & Kluster kontainer men-deploy artefak (Container
Image) yang mewujudkan komponen; server basis data; artefak basis data mewujudkan
objek data; jaringan aman ke BI-FAST. \\
Analisis kesenjangan & Selisih baseline--target: platform yang dibangun,
dimodernisasi, atau dibungkus dari sistem lama. \\
Masukan ke fase berikutnya & Kebutuhan implementasi \& migrasi: work package,
plateau, dan gap deployment (bab berikutnya). \\
\hline
\end{tabularx}
\end{table}

\subsection{Contoh Interaksi Antar-Elemen Teknologi}

Gambar~\ref{fig:technology-contoh} memakai contoh restoran cepat saji: node
\emph{Server Restoran} tersusun atas \emph{OS Restoran} dan mewujudkan
\emph{Layanan Komputasi Restoran} (yang \emph{melayani} komponen \emph{Aplikasi
Kasir}); server itu \emph{men-deploy} artefak \emph{Basis Data Restoran} yang
\emph{mewujudkan} objek data \emph{Data Pesanan}, serta terhubung lewat
\emph{Jaringan Restoran}.

\begin{figure}[htbp]
\centering
\includegraphics[width=\textwidth,height=0.4\textheight,keepaspectratio]{technology-contoh.pdf}
\caption{Contoh interaksi antar-elemen teknologi pada kasus restoran cepat saji.}
\label{fig:technology-contoh}
\end{figure}

\subsection*{Glosarium Mini}
\begin{description}
  \item[Node] Sumber daya komputasi yang men-deploy perangkat lunak.
  \item[Device] Node fisik (perangkat keras).
  \item[System software] Perangkat lunak sistem: OS, runtime, DBMS.
  \item[Technology service] Layanan yang ditawarkan node ke aplikasi.
  \item[Artifact] Wujud fisik data atau komponen yang dapat di-deploy.
  \item[Communication network] Penghubung antar-node.
  \item[Deployment] Penempatan komponen dan artefak pada node.
\end{description}

\section{Kerangka Berpikir}

Alur berpikir bab ini bergerak dari komponen aplikasi menuju infrastruktur
(Gambar~\ref{fig:tech-flow}). Komponen aplikasi menentukan kebutuhan node dan
platform; node men-deploy artefak; artefak mewujudkan komponen dan objek data;
jaringan dan layanan teknologi menopang semuanya; dan pola deployment menata
baseline serta target. Hasilnya menjadi bahan analisis dan kerangka akhir naskah.

\begin{figure}[htbp]
\centering
\scalebox{0.86}{
\begin{tikzpicture}[
  box/.style={rectangle, rounded corners=1.5mm, draw=praditagreen,
    fill=praditagreen!12, thick, align=center, text width=30mm, minimum height=11mm},
  outnode/.style={rectangle, rounded corners=1.5mm, draw=orange!80!black,
    fill=orange!16, thick, align=center, text width=32mm, minimum height=11mm},
  link/.style={-Latex, thick, draw=praditagreen}
]
  \node[box] (comp) {Komponen aplikasi\\(dari Bab 7)};
  \node[box, right=7mm of comp] (node) {Node \&\\perangkat lunak sistem};
  \node[box, right=7mm of node] (art) {Artefak \&\\objek data};
  \node[box, below=8mm of art] (net) {Jaringan \&\\layanan teknologi};
  \node[outnode, left=7mm of net] (deploy) {Pola deployment\\\& kesenjangan};
  \node[outnode, left=7mm of deploy] (an) {Analisis \&\\kerangka akhir};
  \draw[link] (comp) -- (node); \draw[link] (node) -- (art); \draw[link] (art) -- (net);
  \draw[link] (net) -- (deploy); \draw[link] (deploy) -- (an);
  \draw[link, dashed] (an) -- (comp);
\end{tikzpicture}
}
\caption{Alur berpikir lapisan teknologi untuk makalah publikasi.}
\label{fig:tech-flow}
\end{figure}

\section{Konteks Indonesia}

Arsitektur teknologi layanan keuangan Indonesia dibentuk oleh kewajiban
penyelenggaraan sistem elektronik: keandalan, keamanan, pengelolaan pusat data,
dan jejak sistem \cite{pp712019}. Penempatan data (\emph{data residency}) dan
kedaulatan data menjadi pertimbangan penting ketika memilih awan atau pusat data
sendiri. Untuk rel transfer, konektivitas ke BI-FAST menuntut jaringan yang aman
dan andal sesuai standar pembayaran nasional \cite{snapbi}, sedangkan pemrosesan
data pribadi pada node dan artefak harus memenuhi perlindungan data pribadi
\cite{uupdp2022}. Realitas sistem inti lama membuat pendekatan hibrida---kontainer
untuk komponen baru, \emph{mainframe} yang dibungkus adaptor---lebih realistis
daripada penggantian menyeluruh.

\section{Studi Kasus Utama: Arsitektur Teknologi Transfer Bank Wanua}

Komponen aplikasi pada Bab~\ref{ch:application-architecture} di-deploy di atas
sekumpulan node pada Gambar~\ref{fig:technology-transfer}.

\begin{figure}[htbp]
\centering
\includegraphics[width=\textwidth,height=0.62\textheight,keepaspectratio]{technology-transfer.pdf}
\caption{Lanskap teknologi (ArchiMate) API Transfer Dana Bank Wanua: node tersusun
atas perangkat dan perangkat lunak sistem, mewujudkan layanan teknologi, men-deploy
artefak, dan terhubung lewat jaringan.}
\label{fig:technology-transfer}
\end{figure}

\subsection{Cara Membaca Diagram Teknologi}

Gambar~\ref{fig:technology-transfer} memuat tiga \textbf{node} di kiri.
\emph{Kluster Kontainer} \textbf{tersusun atas} (komposisi, ditunjukkan lewat
pelingkupan) \emph{Server Fisik}, \emph{Runtime Kontainer \& OS}, \emph{Platform
Manajemen API}, dan \emph{Message Broker}. \emph{Server Basis Data} tersusun atas
\emph{Sistem Basis Data (PostgreSQL)}. \emph{Sistem Inti (Mainframe)} mewakili
sistem lama dan \textbf{mewujudkan} \emph{Layanan Core Banking} yang melayani
Adaptor Core Banking. Tiap node \textbf{mewujudkan} sebuah \textbf{layanan
teknologi} (Layanan Komputasi Kontainer, Layanan Basis Data, Layanan Core Banking)
dan \textbf{men-deploy} \textbf{artefak} (Container Image, Skema Basis Data, Berkas
Ledger); \emph{Jaringan Internal Bank} mewujudkan \emph{Layanan Jaringan Aman}. Di
bawah, dua \textbf{jaringan komunikasi}---Jaringan Internal Bank dan Koneksi Aman
BI-FAST---\textbf{menghubungkan} node-node tersebut.

\subsection{Penyelarasan Aplikasi, Teknologi, dan Data}

Rantai keterlacakan tertutup ketika lapisan aplikasi, teknologi, dan data
disandingkan (Gambar~\ref{fig:technology-stack}). Diagram ini menyatakan pola
deployment kanonik: \textbf{node men-deploy artefak}, dan \textbf{artefak
mewujudkan} komponen serta data. Di atas, \emph{Container Image}---yang di-deploy
\emph{Kluster Kontainer}---mewujudkan keenam komponen rel transfer. Di bawah,
\emph{Skema Basis Data} dan \emph{Berkas Ledger}---yang di-deploy \emph{Server Basis
Data}---mewujudkan Data Transaksi, Instruksi, dan seterusnya. Node menjalankan
komponen secara \emph{tidak langsung}, melalui artefak yang di-deploy olehnya.

\begin{figure}[htbp]
\centering
\includegraphics[width=\textwidth,height=0.62\textheight,keepaspectratio]{technology-stack.pdf}
\caption{Penyelarasan lintas-lapisan (ArchiMate): komponen aplikasi diwujudkan
oleh artefak (Container Image) yang di-deploy node; objek data diwujudkan oleh
artefak basis data yang di-deploy node basis data.}
\label{fig:technology-stack}
\end{figure}

Dengan ini seluruh rantai lengkap: sasaran motivasi diwujudkan kapabilitas dan
value stream strategi; dijalankan proses, peran, dan layanan bisnis; ditopang
layanan, fungsi, dan komponen aplikasi; yang mengakses objek data; dan pada
akhirnya diwujudkan sebagai artefak yang di-deploy di node pada infrastruktur.

\section{Langkah Analisis dan Artefak Pendukung}

\subsection{Langkah 1: Ambil Komponen Aplikasi dan Objek Data}
Mulai dari komponen aplikasi (Bab~\ref{ch:application-architecture}) dan objek
data (Bab~\ref{ch:data-architecture}) yang perlu di-deploy dan disimpan.

\subsection{Langkah 2: Rancang Node dan Perangkat Lunak Sistem}
Tentukan node yang akan men-deploy artefak, beserta perangkat dan perangkat lunak
sistem yang menyusunnya.

\subsection{Langkah 3: Tentukan Artefak}
Tentukan artefak yang mewujudkan komponen aplikasi dan objek data, serta node
yang men-deploy-nya.

\subsection{Langkah 4: Rancang Jaringan dan Layanan Teknologi}
Tentukan jaringan komunikasi antar-node dan layanan teknologi yang ditawarkan.

\subsection{Langkah 5: Pilih Pola Deployment}
Tetapkan pola deployment (kontainer, awan/pusat data, hibrida) dan alasannya
terhadap kondisi sistem inti dan kepatuhan.

\subsection{Langkah 6: Susun Baseline dan Target}
Bandingkan infrastruktur saat ini dan target, lalu catat kesenjangan.

\subsection{Langkah 7: Kaitkan dengan Regulasi Infrastruktur}
Tautkan penempatan data dan keamanan dengan PP 71 dan konektivitas dengan SNAP BI.

\subsection{Langkah 8: Tulis Narasi Analisis Teknologi}
Ubah lanskap dan matriks menjadi paragraf yang menjelaskan keputusan teknologi dan
implikasinya bagi implementasi.

\begin{table}[htbp]
\centering
\scriptsize
\caption{Artefak pendukung Bab~8 dan fungsinya dalam makalah.}
\label{tab:artefak-bab7-tech}
\begin{tabularx}{\textwidth}{@{}p{0.26\textwidth} p{0.30\textwidth} X@{}}
\hline
\textbf{Artefak} & \textbf{Isi minimum} & \textbf{Dipakai untuk bagian makalah} \\
\hline
Katalog teknologi & Node, perangkat lunak sistem, layanan, artefak & Analisis
domain teknologi. \\
Diagram lanskap teknologi & Node \& penyusunnya, jaringan & Analisis dan lampiran. \\
Diagram penyelarasan lintas-lapisan & Aplikasi--teknologi--data & Keterlacakan
end-to-end \& kerangka akhir. \\
Analisis pola deployment & Pilihan pola + alasan & Diskусi \& implikasi
manajerial. \\
\hline
\end{tabularx}
\end{table}

\section{Integrasi ke Makalah Publikasi}

Lapisan teknologi memperkuat bagian \textbf{Analisis} dan menutup \textbf{kerangka
akhir} naskah. Diagram penyelarasan menjadi bukti bahwa rancangan dapat di-deploy
secara realistis dan patuh. Contoh paragraf:

\begin{quote}
Analisis teknologi menunjukkan bahwa komponen rel transfer paling baik di-deploy
sebagai kontainer pada sebuah kluster, sementara sistem inti lama dibungkus melalui
adaptor---bukan diganti sekaligus. Objek data diwujudkan sebagai artefak basis
data yang berjejak audit, dan konektivitas ke BI-FAST diamankan sesuai kewajiban
penyelenggaraan sistem elektronik.
\end{quote}

\section{Diskusi Kritis, Trade-off, dan Perangkap Umum}

\begin{table}[htbp]
\centering
\scriptsize
\caption{Trade-off dalam arsitektur teknologi.}
\label{tab:tradeoff-bab7-tech}
\begin{tabularx}{\textwidth}{@{}p{0.24\textwidth} X X@{}}
\hline
\textbf{Keputusan} & \textbf{Jika terlalu ringan} & \textbf{Jika terlalu berat} \\
\hline
Awan vs pusat data & Ketergantungan vendor, isu kedaulatan data & Investasi
infrastruktur besar dan lambat. \\
Granularitas node & Satu server besar, sulit diskalakan & Terlalu banyak node,
operasi rumit. \\
Sistem inti lama & Diabaikan, integrasi rapuh & Diganti sekaligus, risiko tinggi. \\
Keamanan \& jejak & Rentan dan tak patuh & Kontrol berlebihan, biaya tinggi. \\
\hline
\end{tabularx}
\end{table}

Perangkap umum: (1) merancang teknologi tanpa menautkannya ke komponen dan data;
(2) mengabaikan sistem inti lama; (3) memilih awan/pusat data tanpa
mempertimbangkan kedaulatan dan kepatuhan; dan (4) melupakan jejak audit serta
keamanan jaringan.

\section{Latihan Terapan Individual}

\subsection*{Latihan~8.1: Lanskap Teknologi}
Untuk organisasi studi kasus Anda, susun lanskap teknologi: node, perangkat lunak
sistem, layanan teknologi, artefak, dan jaringan.

\subsection*{Latihan~8.2: Penyelarasan Lintas-Lapisan}
Buat diagram yang menautkan komponen aplikasi, node/artefak teknologi, dan objek
data untuk satu kapabilitas sasaran.

\subsection*{Latihan~8.3: Pola Deployment dan Baseline--Target}
Pilih pola deployment beserta alasannya, lalu bandingkan infrastruktur baseline
dan target dan catat kesenjangan.

\subsection*{Latihan~8.4: Paragraf Analisis Teknologi}
Tulis satu paragraf analisis teknologi yang menautkan keputusan infrastruktur
dengan kebutuhan aplikasi dan kelayakan implementasi.

\begin{table}[htbp]
\centering
\scriptsize
\caption{Rubrik latihan Bab~8.}
\label{tab:rubrik-bab7-tech}
\begin{tabularx}{\textwidth}{@{}p{0.25\textwidth} X X@{}}
\hline
\textbf{Komponen} & \textbf{Indikator baik} & \textbf{Perlu diperbaiki bila} \\
\hline
Lanskap teknologi & Node, perangkat lunak, artefak jelas \& tertaut & Sekadar
daftar server tanpa relasi. \\
Penyelarasan lintas-lapisan & Menautkan aplikasi--teknologi--data & Teknologi
berdiri sendiri tanpa jejak. \\
Pola deployment & Pilihan beralasan terhadap konteks & Tanpa pertimbangan sistem
lama/kepatuhan. \\
Narasi analisis & Menautkan teknologi ke kebutuhan \& kelayakan & Deskriptif tanpa
argumen. \\
\hline
\end{tabularx}
\end{table}

\section{Ringkasan, Refleksi, dan Jembatan}

\subsection*{Poin Kunci}
\begin{itemize}
  \item Lapisan teknologi memuat node, device, system software, layanan teknologi,
        artefak, dan jaringan.
  \item \textbf{Node men-deploy artefak}; \textbf{artefak mewujudkan} komponen
        aplikasi dan objek data (node menjalankan komponen secara tidak langsung).
        Layanan teknologi \emph{melayani} komponen/fungsi aplikasi.
  \item Deployment menutup rantai keterlacakan dari motivasi sampai infrastruktur.
  \item Pola deployment (kontainer, awan/pusat data, hibrida) menentukan biaya,
        keandalan, dan kepatuhan.
  \item Fase D menghasilkan baseline, target, dan analisis kesenjangan teknologi.
\end{itemize}

\subsection*{Uji Mandiri}
\begin{enumerate}
  \item Jelaskan perbedaan node, device, system software, dan artefak.
  \item Mengapa deployment dinyatakan melalui node yang men-deploy \emph{artefak}
        dan artefak yang mewujudkan komponen serta data---bukan node yang langsung
        ``menjadi'' atau melayani komponen?
  \item Ambil satu komponen aplikasi pada kasus Anda dan sebutkan artefak yang
        mewujudkannya serta node yang men-deploy artefak itu.
\end{enumerate}

\subsection*{Jembatan ke Bab Berikutnya}
Bab berikutnya (Bab~\ref{ch:implementation-migration}) beralih ke \textbf{lapisan
implementasi dan migrasi}: work package, deliverable, plateau, dan gap yang
mengubah arsitektur target menjadi rencana bertahap. Seluruh artefak dari Bab~2 sampai Bab~8---motivasi, strategi, bisnis,
data, aplikasi, dan teknologi---menjadi masukan untuk menyusun peta jalan
implementasi yang realistis.

\section{Bacaan Lanjutan dan Lampiran}

\subsection*{Bacaan Inti}
\begin{itemize}
  \item The Open Group untuk Fase D ADM dan artefak arsitektur teknologi
        \cite{togaf10}.
  \item Newman untuk deployment, kontainer, dan modularisasi
        \cite{newman2021building}.
\end{itemize}

\subsection*{Bacaan untuk Memperdalam}
\begin{itemize}
  \item Untuk hubungan teknologi dengan operating model dan fondasi eksekusi
        \cite{ross2006enterprise}.
  \item Untuk konteks Indonesia: penyelenggaraan sistem elektronik dan pusat data
        \cite{pp712019}, standar API pembayaran \cite{snapbi}, dan perlindungan
        data pribadi \cite{uupdp2022}.
\end{itemize}

\subsection*{Lampiran Kerja Bab~8}
\begin{description}
  \item[Katalog teknologi] Daftar node, perangkat lunak sistem, layanan, artefak.
  \item[Diagram lanskap teknologi] Node dan penyusunnya.
  \item[Diagram penyelarasan] Aplikasi--teknologi--data.
  \item[Analisis deployment] Pola deployment dan kesenjangan baseline--target.
\end{description}
\setlength{\emergencystretch}{0pt}
